% Part: first-order-logic
% Chapter: axiomatic-deduction
% Section: provability-propositional

\documentclass[../../../include/open-logic-section]{subfiles}

\begin{document}

\iftag{FOL}
      {\olfileid{fol}{axd}{ppr}}
      {\olfileid{pl}{axd}{ppr}}

\olsection{\usetoken{S}{derivability} మరియు ప్రతిజ్ఞాత్మక సంయోజకాలు}

\begin{explain}
  $!A \land !B \Proves !A$, $!A, !A \lif !B \Proves !B$ (మోడస్
  పోనెన్స్) వంటి, ప్రతిజ్ఞాత్మక సంయోజకాలకు సంబంధించిన కొన్ని ప్రాథమిక
  వాస్తవాలను స్థాపించడానికి స్వీకృతాధారిత నిగమనంలోని
  !!{derivability} సంబంధం~$\Proves$ తగినంత బలమైనదని నిరూపిస్తాం.
  సంపూర్ణతా సిద్ధాంత నిరూపణకు ఈ వాస్తవాలు అవసరం.
\end{explain}

\begin{prop}\ollabel{prop:provability-land}
  \begin{enumerate}
  \item \ollabel{prop:provability-land-left} $!A \land !B \Proves
    !A$, $!A \land !B \Proves !B$ రెండూ.
  \item \ollabel{prop:provability-land-right} $!A, !B \Proves !A \land !B$.
  \end{enumerate}
\end{prop}

\begin{proof}
  \begin{enumerate}
    \item \olref[prp]{ax:land1}, \texttt{ax:land2}ల నుంచి మోడస్
      పోనెన్స్ ద్వారా.
    \sourcecorrection{OLTEAXD-007}{సంయోగపు ఎడమ, కుడి భాగాలను పొందే మొదటి
    నిరూపణలో ఆంగ్ల మూలం రెండవ ప్రక్షేపణకూ \olref[prp]{ax:land1}ను
    సూచించింది. రెండవ పర్యవసానానికి అవసరమైనది \texttt{ax:land2} కనుక
    ఆ సూచనను సరిచేశాం.}
  \item \olref[prp]{ax:land3} నుంచి మోడస్ పోనెన్స్‌ను రెండుసార్లు
    వర్తింపజేయడం ద్వారా.
  \end{enumerate}
\end{proof}

\begin{prop}\ollabel{prop:provability-lor}
  \begin{enumerate}
  \item $!A \lor !B, \lnot !A, \lnot !B$ విరుద్ధం.
  \item $!A \Proves !A \lor !B$, $!B \Proves !A \lor !B$ రెండూ.
  \end{enumerate}
\end{prop}

\begin{proof}
  \begin{enumerate}
  \item \texttt{ax:lnot2} నుంచి $\Proves \lnot !A \lif (!A
    \lif \lfalse)$, $\Proves \lnot !B \lif (!B \lif \lfalse)$
    లభిస్తాయి.
    \sourcecorrection{OLTEAXD-008}{ఆంగ్ల మూలం $\lnot !A \lif
    (!A \lif \lfalse)$ను పొందడానికి \olref[prp]{ax:lnot1}ను
    సూచించింది. ఆ రూపం \texttt{ax:lnot2} యొక్క రూపం కనుక సూచనను
    సరిచేశాం.}
    కాబట్టి నిగమన సిద్ధాంతం వల్ల $\{\lnot !A\} \Proves
    !A \lif \lfalse$, $\{\lnot !B\} \Proves !B \lif \lfalse$.
    \olref[prp]{ax:lor3} నుంచి $\{\lnot !A, \lnot !B\} \Proves
    (!A \lor !B) \lif \lfalse$ లభిస్తుంది. నిగమన సిద్ధాంతం వల్ల,
    $\{!A \lor !B, \lnot !A, \lnot !B\} \Proves \lfalse$.
  \item \olref[prp]{ax:lor1}, \olref[prp]{ax:lor2}ల నుంచి మోడస్
    పోనెన్స్ ద్వారా.
  \end{enumerate}
\end{proof}

\begin{prop}\ollabel{prop:provability-lif}
  \begin{enumerate}
  \item \ollabel{prop:provability-lif-left} $!A, !A \lif !B \Proves !B$.
  \item \ollabel{prop:provability-lif-right}
    $\lnot !A \Proves !A \lif !B$, $!B \Proves !A \lif !B$ రెండూ.
  \end{enumerate}
\end{prop}

\begin{proof}
  \begin{enumerate}
  \item మనం ఇలా !!{derive} చేయవచ్చు:
    \begin{derivation}
      1. & $!A$ & \Hyp\\
      2. & $!A \lif !B$ & \Hyp\\
      3. & $!B$ & 1, 2, \MP
    \end{derivation}

  \item వరుసగా \olref[prp]{ax:lnot2}, \olref[prp]{ax:lif1}, నిగమన
    సిద్ధాంతం ద్వారా.
  \end{enumerate}
\end{proof}

\end{document}
